\chapter{Physics Engine Comparison}
\label{ch:engine-comparison}

An engine comparison needs a specified question, model and accuracy target.
Solvers may disagree because their contact laws differ, an import changed the
inertia, or an integration is inaccurate. Conversely, agreement can preserve a
shared modeling mistake. This chapter gives a comparison protocol and executable
error metric; it presents no measured cross-engine speed ranking or energy-drift
curves.

\section{Establish Model Equivalence First}

Begin with the contact-free cylinder pendulum of Chapter~\ref{ch:platform}.
Match the fixed base, joint frame and axis, COM inertia, gravity, damping,
initial state, force inputs and duration. Verify loaded parameters rather than
assuming identical XML implies identical interpreted physics. At several states,
compare mass, gravitational generalized force and passive torque against the
analytic values $I_P$, $mgL\sin q/2$ and $-b\dot q$.

Next vary solver tolerances or internal step sizes and compare trajectories at
common timestamps. An adaptive solver's output interval is not its internal
step size. For fixed-step methods, record substeps and control-update timing.
Interpolation introduces its own error, especially across impacts. Keep
event-time and pre/post-impact comparisons separate from smooth-trajectory
interpolation.

\section{Measure the Right Error}

For the unforced pendulum, define the energy-work residual
\begin{equation}
 r_E(t)=E(t)-E(0)+\int_0^t b\dot q(s)^2\,ds.
\end{equation}
With applied generalized torque $\tau$, subtract
$\int_0^t\tau\dot q\,ds$ as well. Numerical integration of work contributes
to residual error. With impacts or other dissipative forces, include the
corresponding loss or jump law. A percentage relative to $E(0)$ is undefined
when $E(0)=0$ and sensitive to the arbitrary potential-energy reference; report
joules or a declared positive characteristic scale instead.

Energy conservation does not guarantee phase accuracy. A trajectory can stay
near the correct energy surface while reaching the wrong position at a specified
time. For aligned Euclidean states $x_i$ and reference states $x_i^*$, choose
positive scales $s_j$ with each coordinate's units and define
\begin{equation}
 e_{ij}=\frac{x_{ij}-x_{ij}^*}{s_j},\qquad
 E_{\rm RMS}=\left[\frac{1}{T}\int_{t_0}^{t_1}\|e(t)\|^2dt\right]^{1/2},
 \quad T=t_1-t_0.
\end{equation}
The scales make mixed angle/rate or position/velocity comparisons dimensionless
and declare the relative importance assigned to each component. This code uses
trapezoidal quadrature and reports the largest sampled norm, not a certified
continuous-time maximum.

\begin{lstlisting}[language=Python, caption={Time-Weighted Errors on an Aligned Euclidean Grid}]
import numpy as np
from scipy.integrate import trapezoid


def trajectory_errors(times: np.ndarray, actual: np.ndarray,
                      reference: np.ndarray,
                      scales: np.ndarray) -> dict[str, float]:
    """Return dimensionless RMS and sampled maximum normalized state errors."""
    times = np.asarray(times, dtype=float)
    actual = np.asarray(actual, dtype=float)
    reference = np.asarray(reference, dtype=float)
    scales = np.asarray(scales, dtype=float)
    if (times.ndim != 1 or times.size < 2
            or not np.all(np.isfinite(times))):
        raise ValueError("times must be finite 1D with >= 2 samples")
    if np.any(np.diff(times) <= 0):
        raise ValueError("times must strictly increase")
    if (actual.ndim != 2 or actual.shape != reference.shape
            or actual.shape[0] != times.size or actual.shape[1] == 0):
        raise ValueError("states must have matching nonempty (N, d) shapes")
    if not np.all(np.isfinite(actual)) or not np.all(np.isfinite(reference)):
        raise ValueError("states must be finite")
    if (scales.shape != (actual.shape[1],)
            or not np.all(np.isfinite(scales)) or np.any(scales <= 0)):
        raise ValueError("one finite positive scale is required per coordinate")
    squared = np.sum(((actual - reference) / scales)**2, axis=1)
    mean_square = trapezoid(squared, times) / (times[-1] - times[0])
    return {"rms": float(np.sqrt(mean_square)),
            "max": float(np.sqrt(np.max(squared)))}
\end{lstlisting}

For $t=(0,1,3)$ and scalar normalized errors $(0,1,3)$, the trapezoidal integral
of squared error is $10.5$, giving $\sqrt{3.5}$. The unweighted node RMS is
$\sqrt{10/3}$; neither equals the exact continuous RMS $\sqrt{3}$ for
$e(t)=t$ on this coarse grid. Refinement matters.

The caller must first construct physically appropriate differences. Subtracting
quaternion components is not an orientation error: $q$ and $-q$ encode the same
rotation. Use a declared relative-rotation logarithm and tangent-frame convention,
with care near angle $\pi$. Periodic angles require a wrapping or unwrapping
convention. Two simulations' difference is an agreement metric; calling it
accuracy requires an independently justified reference, such as an analytic
limit or a documented convergence study.

\section{Contact Law, Solver and Integrator Are Separate Choices}

The contact model specifies forces or impulses; the solver computes their
coupled solution; the integrator advances state. Collision geometry and contact
detection add further approximations. An engine name does not identify all
these choices.

\begin{description}[style=nextline]
 \item[MuJoCo.] Its soft-contact formulation uses convex optimization and
 relaxes strict complementarity. Contact parameters describe the selected
 compliant behavior; integration method and step settings are separate choices.
 A nominal penetration value is neither automatically an error nor evidence
 of validated tissue deformation \cite{MuJoCoComputationReview}.
 \item[Drake.] Hydroelastic contact uses pressure fields and contact surfaces
 to model compliant interaction while bodies retain rigid-body motion. It does
 not resolve arbitrary elastic waves inside bodies. Geometry representations,
 compliant/rigid pairing and point-contact fallback settings affect which
 contacts are supported \cite{DrakeHydroelasticReview}.
 \item[Pinocchio.] Constrained and impulse dynamics are available; saying
 ``no contact dynamics'' is incorrect. Supplying constraints differs from
 implementing collision detection, friction and time integration together.
 Identify the additional model and algorithms actually used
 \cite{PinocchioAlgorithmsReview}.
 \item[OpenSim.] Behavior depends on the selected force component. For example,
 the documented HuntCrossleyForce uses Hertz-based interactions between spheres
 and half-spaces, with material and friction parameters. That component's
 geometry restrictions do not describe every OpenSim contact component
 \cite{OpenSimHuntCrossleyReview}.
\end{description}

Changing contact stiffness changes the physical model and can also make its
equations numerically harder to resolve. Report both effects. In golf,
foot-ground loading, hand-club constraints and club-ball impact can require
different resolutions. Agreement on a smooth pendulum does not validate these
contact regimes.

\section{A Reproducible Comparison Record}

For each run record model and parameter hashes, engine and wrapper versions,
gravity and force conventions, solver and integrator options, internal steps or
tolerances, output and command timing, and initial state. Record computer,
thread configuration, precision, initialization/warm-up policy and repeated
timing trials. Separate loading, compilation, stepping and output costs. Report
timing variability and failed runs, not just the fastest successful sample.

Compare cost at a declared error tolerance as well as error at a declared cost.
Equal nominal timesteps need not yield equal accuracy. Start with analytic
limits, then smooth multibody cases, and then the required contact regime.
With experimental data, propagate measurement and parameter uncertainty: a
converged solver cannot compensate for an incorrect joint axis or unmeasured
external force. These checks support conclusions about a specified model and
observable, not an unrestricted winner among engines.

\section*{Exercises}
\begin{enumerate}
 \item \textbf{Pendulum Refinement.} Use the same physical cylinder in two
 available implementations. Report parameter agreement, phase error, work
 residual and runtime across refinements; identify the reference's own error.
 \item \textbf{Nonuniform Sampling.} Reproduce the three-node example, refine
 the grid, and distinguish quadrature error from simulation error.
 \item \textbf{Contact Reference.} Compare a bouncing-ball model with the
 ballistic flight and restitution law it assumes. Ideal elastic impact is a
 suitable reference only for that limit, not automatically for dissipative
 compliant contact.
 \item \textbf{Swing Observable.} Choose clubhead velocity, joint torque or
 contact impulse. State units, alignment, uncertainty and model assumptions
 before choosing a metric.
\end{enumerate}
